% Additional lessons: heaps, search, hashing, and house robber DP.

\patternintoc{Heaps and Streaming Data}
\problemstart{Find Median from Data Stream}{295}{Hard}{Two balanced heaps}{Priority queues}

\subsection{The Problem}
Design a class that accepts integers one at a time and returns the median of
all values seen so far. An odd-sized stream has one middle value; an even-sized
stream uses the average of the two middle values.

\subsection{Two-Heap Model}
A sorted vector makes median lookup easy but insertion costs \(O(n)\). Instead,
split the stream into two halves:
\begin{itemize}
  \item \texttt{lower}: a max heap containing the smaller half;
  \item \texttt{upper}: a min heap containing the larger half.
\end{itemize}
Keep \(|\texttt{lower}|=|\texttt{upper}|\) or one extra item in
\texttt{lower}, and keep every lower value no greater than every upper value.

\begin{invariantbox}{the heap tops are always the middle boundary}
After every insertion, all lower-heap values are no greater than all
upper-heap values, and the lower heap has either the same size or one extra
value. These two facts place the median at \texttt{lower.top()} or between the
two heap tops.
\end{invariantbox}

\begin{visualbox}{The two top values give the median}
\centering
\begin{tikzpicture}[>={Stealth[length=2.2mm,width=1.55mm]},font=\small]
  \node[dsnode,fill=softblue,minimum width=42mm] (lo)
    {smaller half\\largest here \(=3\)\\values \(\{1,2,3\}\)};
  \node[dsnode,fill=softteal,minimum width=42mm,right=18mm of lo] (hi)
    {larger half\\smallest here \(=4\)\\values \(\{4,5,8\}\)};
  \draw[flowarrow] (lo)--node[edgelabel,above]{middle split} (hi);
  \node[below=8mm of $(lo)!0.5!(hi)$,draw=orange,fill=softorange,rounded corners]
    {median \(=(3+4)/2=3.5\)};
\end{tikzpicture}
\end{visualbox}

\subsection{Adding a Number}
Push the new value into \texttt{lower}. Move the largest lower value to
\texttt{upper}; this restores ordering. If \texttt{upper} becomes larger, move
its smallest value back to \texttt{lower}. This restores the size rule.

\subsection{Pseudocode}
\begin{lstlisting}[style=pseudocode]
addNum(value):
    push value into lower max heap
    move lower.maximum into upper min heap
    if upper has more values than lower:
        move upper.minimum back into lower

findMedian():
    if lower has one extra value: return lower.maximum
    return average(lower.maximum, upper.minimum)
\end{lstlisting}

\subsection{C++ Implementation}
\begin{lstlisting}[style=compactcode,caption={Complete C++ data structure: MedianFinder}]
#include <functional>
#include <queue>
#include <vector>

using namespace std;

class MedianFinder {
    priority_queue<int> lower;
    priority_queue<int, vector<int>, greater<int>> upper;

public:
    void addNum(int num) {
        lower.push(num);
        upper.push(lower.top());
        lower.pop();

        if (upper.size() > lower.size()) {
            lower.push(upper.top());
            upper.pop();
        }
    }

    double findMedian() const {
        if (lower.size() > upper.size()) return lower.top();
        return (static_cast<double>(lower.top()) + upper.top()) / 2.0;
    }
};
\end{lstlisting}

\subsection{Why It Works and Complexity}
Moving the maximum lower value into \texttt{upper} ensures all lower values
remain no greater than all upper values. Rebalancing changes sizes by one, so
the middle value or values remain at the heap tops. Insertion costs
\(O(\log n)\), median lookup costs \(O(1)\), and stored space is \(O(n)\).

\begin{revisionbox}{Common mistakes}
Use a max heap for the lower half and a min heap for the upper half. Convert to
\texttt{double} before averaging to avoid integer truncation and overflow.
\end{revisionbox}

\subsection{From a Sorted Vector to Two Heaps}
Keeping every value in a sorted vector makes the median easy to read, but an
insertion may shift \(O(n)\) values. Sorting from scratch after every insertion
is even more expensive. Two heaps keep only the ordering information needed
for the middle: the maximum of the lower half and the minimum of the upper
half.

The keeped invariants are:
\begin{enumerate}
  \item every value in \texttt{lower} is at most every value in
  \texttt{upper};
  \item the heaps have equal sizes or \texttt{lower} has one extra value.
\end{enumerate}

\subsection{Insertion Dry Run}
\begin{dryrunbox}{Adding \(5,2,8,3\)}
\begin{center}
\begin{tabular}{@{}c c c c@{}}
\toprule
Inserted & Lower max heap & Upper min heap & Median \\
\midrule
5 & \(\{5\}\) & \(\{\}\) & 5 \\
2 & \(\{2\}\) & \(\{5\}\) & 3.5 \\
8 & \(\{5,2\}\) & \(\{8\}\) & 5 \\
3 & \(\{3,2\}\) & \(\{5,8\}\) & 4 \\
\bottomrule
\end{tabular}
\end{center}
The braces show heap contents conceptually; only each heap's top is ordered
for direct access.
\end{dryrunbox}

\begin{visualbox}{The top of each heap is one middle value}
\centering
\begin{tikzpicture}[font=\small,>={Stealth[length=2.2mm,width=1.55mm]}]
  \node[trienode,active] (l0) at (0,1.2) {3};
  \node[trienode,below=10mm of l0] (l1) {2};
  \draw[-,thick,navy] (l0)--(l1);
  \node[above=5mm of l0,text=orange,font=\bfseries]{smaller half};
  \node[below=5mm of l1,font=\scriptsize]{top = largest in this half};

  \node[trienode,visited] (u0) at (5.2,1.2) {5};
  \node[trienode,below=10mm of u0] (u1) {8};
  \draw[-,thick,navy] (u0)--(u1);
  \node[above=5mm of u0,text=teal,font=\bfseries]{larger half};
  \node[below=5mm of u1,font=\scriptsize]{top = smallest in this half};

  \draw[<->,very thick,navy] (l0.east)--node[edgelabel,above]{middle split} (u0.west);
  \node[below=17mm of $(l0)!0.5!(u0)$,draw=orange,fill=softorange,
    rounded corners]{median \(=(3+5)/2=4\)};
\end{tikzpicture}
\end{visualbox}

\subsection{Code Walkthrough}
Every new value first enters \texttt{lower}. Moving \texttt{lower.top()} to
\texttt{upper} guarantees the ordering boundary. If upper now has more
values, moving its minimum back restores the size rule. The median is
then one lower top for an odd count or the average of both tops for an even
count.

\subsection{Edge Cases, Overflow, and Follow-ups}
The problem calls \texttt{findMedian} only after at least one insertion.
Duplicate and negative values require no special handling. Cast before adding
the two middle values; otherwise two large integers can overflow before being
converted. A harder follow-up supports deletions or a sliding window, usually
with lazy deletion maps because priority queues cannot erase arbitrary values
efficiently.

\begin{revisionbox}{How to explain it simply}
I split the stream around the median. A max heap exposes the largest lower
value and a min heap exposes the smallest upper value. After every insertion I
restore ordering and a one-element size difference, giving \(O(\log n)\)
updates and \(O(1)\) median queries.
\end{revisionbox}
\chaptercheckpoint
  {Values arrive online and the middle must remain quickly accessible.}
  {A lower max heap and upper min heap with ordering and size invariants.}
  {Using two heaps of the same type or averaging before converting to double.}
  {$O(\log n)$ insertion, $O(1)$ query, and $O(n)$ storage.}
\judgeclassnote


\patternintoc{Binary Search}
\problemstart{Find Minimum in Rotated Sorted Array}{153}{Medium}{Binary search}{Array}

\subsection{The Problem}
A strictly increasing array has been rotated. Return its minimum value in
\(O(\log n)\) time. All values are distinct.

\begin{conceptbox}{Example}
\texttt{[4,5,6,7,0,1,2]} has minimum \(0\). An unrotated array such as
\texttt{[1,2,3]} has minimum \(1\).
\end{conceptbox}

\subsection{Binary-Search Decision}
Compare the midpoint with the right endpoint:
\begin{itemize}
  \item If \(nums[mid]>nums[right]\), the rotation point is strictly right of
  \(mid\), so set \(left=mid+1\).
  \item Otherwise, \(mid\) may itself be the minimum, so set \(right=mid\).
\end{itemize}

\begin{invariantbox}{the closed interval always contains the rotation minimum}
The comparison with the right endpoint identifies a half that cannot contain
the minimum. The algorithm removes only that half, retaining \texttt{mid}
when it may itself be the answer, until both boundaries meet at the minimum.
\end{invariantbox}

\begin{visualbox}{Because \(7>2\), the minimum must be to the right}
\centering
\begin{tikzpicture}[>={Stealth[length=2.2mm,width=1.55mm]},font=\small]
  \foreach \x/\v in {0/4,1/5,2/6,3/7,4/0,5/1,6/2} {
    \node[arraycell] (a\x) at (1.1*\x,0) {\v};
  }
  \node[above=6mm of a0,text=teal] {left};
  \node[above=6mm of a3,text=orange] {mid};
  \node[above=6mm of a6,text=navy] {right};
  \draw[flowarrow,orange] (a3.south) to[bend right=25]
    node[edgelabel,below]{keep the right half} (a4.south);
\end{tikzpicture}
\end{visualbox}

\subsection{Pseudocode}
\begin{lstlisting}[style=pseudocode]
left = 0
right = nums.length - 1
while left < right:
    middle = left + (right - left) / 2
    if nums[middle] > nums[right]:
        left = middle + 1
    else:
        right = middle
return nums[left]
\end{lstlisting}

\subsection{C++ Implementation}
\begin{lstlisting}[style=compactcode,caption={Complete C++ solution: Find Minimum in Rotated Sorted Array}]
#include <vector>

using namespace std;

class Solution {
public:
    int findMin(vector<int>& nums) {
        int left = 0;
        int right = static_cast<int>(nums.size()) - 1;
        while (left < right) {
            const int middle = left + (right - left) / 2;
            if (nums[middle] > nums[right]) left = middle + 1;
            else right = middle;
        }
        return nums[left];
    }
};
\end{lstlisting}

\subsection{Why It Works, Complexity, and Common Mistakes}
The search interval always contains the minimum. Each comparison safely
discards one half while retaining \(mid\) when it may be the answer. Interval
size halves each iteration: \(O(\log n)\) time and \(O(1)\) space. Use
\(right=mid\), not \(mid-1\), in the second branch.

\subsection{From Linear Scan to Binary Search}
A linear scan can return the smallest value in \(O(n)\), but the rotated
sorted structure allows a logarithmic decision. Comparing the midpoint with
the right endpoint reveals whether the midpoint lies in the high rotated
portion or in the low portion containing the minimum.

\subsection{Boundary Dry Run}
\begin{dryrunbox}{Searching \texttt{[4,5,6,7,0,1,2]}}
\begin{center}
\begin{tabular}{@{}c c c c c@{}}
\toprule
Left & Mid & Right & Comparison & New interval \\
\midrule
0 (4) & 3 (7) & 6 (2) & \(7>2\) & \([4,6]\) \\
4 (0) & 5 (1) & 6 (2) & \(1\le2\) & \([4,5]\) \\
4 (0) & 4 (0) & 5 (1) & \(0\le1\) & \([4,4]\) \\
\bottomrule
\end{tabular}
\end{center}
The loop stops when both boundaries identify the minimum value 0.
\end{dryrunbox}

\subsection{Why Compare with the Right Endpoint?}
If \(nums[mid]>nums[right]\), the sorted order breaks somewhere strictly to
the right of mid, so mid cannot be the minimum. Otherwise mid belongs to the
low sorted portion and may itself be the minimum, which is why the assignment
is \texttt{right = middle} rather than \texttt{middle - 1}.

\subsection{Edge Cases and Variations}
A one-element or unrotated array works without a special branch. The proof
uses distinct values. If duplicates are allowed, equality may not identify a
safe half; a common adaptation decrements the right boundary on equality,
which can degrade to \(O(n)\). Compute the midpoint as
\texttt{left + (right-left)/2} to avoid integer overflow.

\begin{revisionbox}{Practice checkpoint}
Trace the algorithm on \texttt{[2,3,4,5,1]} and \texttt{[1,2,3,4]}. At every
step state why the discarded half cannot contain the minimum.
\end{revisionbox}
\chaptercheckpoint
  {A sorted array has one rotation break and logarithmic time is required.}
  {A closed binary-search interval containing the minimum.}
  {Using \texttt{right=mid-1} when the midpoint may be the answer.}
  {$O(\log n)$ time and $O(1)$ space for distinct values.}
\judgeclassnote


\patternintoc{Undirected Graphs}
\problemstart{Graph Valid Tree}{261}{Medium}{DFS cycle and connectivity check}{Adjacency list}

\subsection{The Problem}
Given \(n\) labeled vertices and undirected edges, return whether the graph is
a tree. A tree must be connected and contain no cycle.

\subsection{Two Conditions}
An undirected tree with \(n\) vertices has exactly \(n-1\) edges. That count is
a quick rejection. Then DFS from vertex zero, ignoring the edge back to the
parent. Reaching any other visited neighbor proves a cycle. Finally, every
vertex must have been reached.

\begin{invariantbox}{visited nodes form the explored part of one rooted traversal}
Every recursive frame records the edge used to enter its node as the parent
edge. That one reverse edge is expected in an undirected graph and is skipped;
any other visited neighbor closes a cycle. After DFS, full visitation proves
that no disconnected component was omitted.
\end{invariantbox}

\begin{visualbox}{A tree is connected and has no loop}
\centering
\begin{tikzpicture}[>={Stealth[length=2.2mm,width=1.55mm]},font=\small]
  \node[trienode,visited] (a) {0}; \node[trienode,visited,below left=12mm and 12mm of a] (b) {1};
  \node[trienode,visited,below right=12mm and 12mm of a] (c) {2};
  \draw[-,teal,thick] (a)--(b); \draw[-,teal,thick] (a)--(c);
  \node[below=7mm of $(b)!0.5!(c)$,text=teal] {valid tree};
  \node[trienode,right=42mm of a] (x) {0}; \node[trienode,below left=12mm and 10mm of x] (y) {1};
  \node[trienode,below right=12mm and 10mm of x] (z) {2};
  \draw[-,orange,thick] (x)--(y)--(z)--(x);
  \node[below=7mm of $(y)!0.5!(z)$,text=orange] {loop: not a tree};
\end{tikzpicture}
\end{visualbox}

\begin{visualbox}{Neighbor lists, visited marks, and the path being checked}
\centering
\begin{tikzpicture}[font=\small,>={Stealth[length=2.2mm,width=1.55mm]}]
  \node[dsnode,align=left,fill=softteal] (adj) {
    \textbf{neighbors}\\
    \texttt{0: [1,2,3]}\\
    \texttt{1: [0,4]}\\
    \texttt{2: [0]}\\
    \texttt{3: [0]}\\
    \texttt{4: [1]}};
  \node[draw=navy,fill=softblue,minimum width=42mm,minimum height=8mm,
    right=27mm of adj] (f0) {visit 0; no parent};
  \node[draw=navy,fill=softblue,minimum width=42mm,minimum height=8mm,
    above=0mm of f0] (f1) {visit 1; came from 0};
  \node[draw=orange,fill=softorange,minimum width=42mm,minimum height=8mm,
    above=0mm of f1] (f2) {visit 4; came from 1};
  \node[right=6mm of f2,text=orange,font=\bfseries,align=left] {working\\now};
  \draw[flowarrow,teal] (adj.east)--(f0.west);
  \foreach \x/\v in {0/T,1/T,2/F,3/F,4/T} {
    \node[arraycell,minimum width=10mm] (seen\x) at ({3.0+1.1*\x},-1.45) {\v};
    \node[below=1mm of seen\x,font=\scriptsize] {\x};
  }
  \node[below=7mm of seen2,text=navy,font=\bfseries] {nodes already seen};
\end{tikzpicture}

\smallskip
The list shows connected nodes, the stack remembers how we arrived, and the
marks show which nodes were already visited.
\end{visualbox}

\subsection{Pseudocode}
\begin{lstlisting}[style=pseudocode]
if edge count is not n - 1: return false
build an undirected adjacency list
dfs(node, parent):
    if node is already visited: return false
    mark node visited
    for neighbor in graph[node]:
        skip parent
        if dfs(neighbor, node) is false: return false
    return true
run dfs from node 0
return dfs succeeded and every node was visited
\end{lstlisting}

\subsection{C++ Implementation}
\begin{lstlisting}[style=compactcode,caption={Complete C++ solution: Graph Valid Tree}]
#include <vector>

using namespace std;

class Solution {
    bool dfs(int node, int parent,
             const vector<vector<int>>& graph,
             vector<bool>& visited) {
        if (visited[node]) return false;
        visited[node] = true;
        for (const int neighbor : graph[node]) {
            if (neighbor == parent) continue;
            if (!dfs(neighbor, node, graph, visited)) return false;
        }
        return true;
    }

public:
    bool validTree(int n, vector<vector<int>>& edges) {
        if (static_cast<int>(edges.size()) != n - 1) return false;
        vector<vector<int>> graph(n);
        for (const auto& edge : edges) {
            graph[edge[0]].push_back(edge[1]);
            graph[edge[1]].push_back(edge[0]);
        }
        vector<bool> visited(n, false);
        if (!dfs(0, -1, graph, visited)) return false;
        for (const bool reached : visited) if (!reached) return false;
        return true;
    }
};
\end{lstlisting}

\subsection{Why It Works and Complexity}
The parent check prevents treating the same undirected edge as a cycle. Any
other repeated vertex closes a cycle. The final visited check proves
connectivity. DFS and adjacency construction take \(O(V+E)\) time and space.

\begin{revisionbox}{How to explain it in an interview}
A tree is connected and acyclic. I first verify the needed \(n-1\) edge
count, then DFS while passing the parent. A visited non-parent neighbor is a
cycle; any unvisited vertex afterward means the graph is disconnected.
\end{revisionbox}

\subsection{Why Both Conditions Matter}
A connected graph may still contain a cycle, and an acyclic graph may be a
disconnected forest. So checking only one property is inenough.
The \(n-1\) edge count is a fast needed condition, while DFS verifies the
actual structure.

\subsection{DFS Dry Run}
\begin{dryrunbox}{Edges \texttt{[[0,1],[0,2],[0,3],[1,4]]}}
\begin{center}
\begin{tabularx}{\linewidth}{@{}c c X@{}}
\toprule
Current & Parent & Action \\
\midrule
0 & -1 & mark 0; recurse to 1, 2, and 3 \\
1 & 0 & ignore edge back to 0; recurse to 4 \\
4 & 1 & ignore edge back to 1; return true \\
2, 3 & 0 & each is new and has only its parent edge \\
finish & -- & all five vertices are visited; return true \\
\bottomrule
\end{tabularx}
\end{center}
An extra edge \([1,2]\) would lead from one DFS branch to an already visited
non-parent vertex and prove a cycle.
\end{dryrunbox}

\subsection{Code Walkthrough}
The edge-count guard rejects obvious non-trees before allocating the graph.
Each undirected edge is stored in both adjacency lists. DFS marks a node,
skips exactly the edge to its parent, and rejects any other repeated node. The
final loop confirms that vertex zero's traversal reached every component.

\subsection{Alternative and Edge Cases}
Union find is a high-value alternative: union every edge and reject an edge
whose endpoints already share a representative, then confirm \(n-1\) edges.
For \(n=1\), an empty edge list is a valid tree. Self-loops and parallel edges
create cycles. The recursive DFS may use \(O(n)\) stack space on a chain.

\begin{revisionbox}{Practice checkpoint}
Explain why the parent edge is skipped but any other visited neighbor is a
cycle. Then outline the union-find version and compare its state with DFS.
\end{revisionbox}
\chaptercheckpoint
  {An undirected graph must be both connected and acyclic.}
  {An adjacency list, visited array, DFS parent, and the \(n-1\) edge test.}
  {Treating the parent edge as a cycle or checking only connectivity.}
  {$O(V+E)$ time and $O(V+E)$ space.}
\judgeclassnote


\patternintoc{Arrays and Hashing}
\problemstart{Group Anagrams}{49}{Medium}{Standard-key hashing}{Hash map}

\subsection{The Problem}
Group strings that contain the same letters with the same frequencies. The
groups and strings within them may be returned in any order.

\subsection{Sorted Standard Key}
Anagrams become identical after sorting their characters. Use the sorted word
as a hash-map key and append the original word to that bucket.

\begin{invariantbox}{each processed word is stored in its unique standard bucket}
After processing a prefix of the input, every word from that prefix appears in
the bucket indexed by its sorted characters. Two words share a bucket exactly
when their character multisets match, so no later comparison between words is
needed.
\end{invariantbox}

\begin{visualbox}{Sort each word to find its group}
\centering
\begin{tikzpicture}[>={Stealth[length=2.2mm,width=1.55mm]},font=\small]
  \node[dsnode] (a) {\texttt{eat}}; \node[dsnode,below=5mm of a] (b) {\texttt{tea}};
  \node[dsnode,below=5mm of b] (c) {\texttt{ate}};
  \node[dsnode,right=25mm of b,fill=softorange] (key) {sorted letters\\\texttt{aet}};
  \draw[flowarrow] (a)--(key); \draw[flowarrow] (b)--(key); \draw[flowarrow] (c)--(key);
  \node[dsnode,right=22mm of key,fill=softteal] (bucket)
    {one group\\\texttt{[eat,tea,ate]}};
  \draw[flowarrow,teal] (key)--(bucket);
\end{tikzpicture}
\end{visualbox}

\subsection{Pseudocode}
\begin{lstlisting}[style=pseudocode]
groups = empty map from sorted key to words
for word in strings:
    key = characters of word sorted
    append word to groups[key]
return all map buckets
\end{lstlisting}

\subsection{C++ Implementation}
\begin{lstlisting}[style=compactcode,caption={Complete C++ solution: Group Anagrams}]
#include <algorithm>
#include <string>
#include <unordered_map>
#include <utility>
#include <vector>

using namespace std;

class Solution {
public:
    vector<vector<string>>
    groupAnagrams(vector<string>& strings) {
        unordered_map<string, vector<string>> groups;
        for (const string& word : strings) {
            string key = word;
            sort(key.begin(), key.end());
            groups[key].push_back(word);
        }

        vector<vector<string>> result;
        result.reserve(groups.size());
        for (auto& entry : groups) result.push_back(move(entry.second));
        return result;
    }
};
\end{lstlisting}

\subsection{Complexity and Alternatives}
For \(n\) strings of maximum length \(k\), sorting keys costs
\(O(nk\log k)\); stored keys and output use \(O(nk)\) space. A 26-count
frequency signature avoids sorting and gives \(O(nk)\) time, but the sorted
key is simpler and easy to explain.

\begin{revisionbox}{Common mistakes}
Append the original word, not the sorted key. Empty strings correctly share an
empty key. Do not assume the order of buckets from an unordered map.
\end{revisionbox}

\subsection{From Pairwise Comparison to a Standard Form}
A direct method can compare every pair of strings by character frequencies,
which costs roughly \(O(n^2k)\). The repeated work is deciding the same
anagram identity again and again. A standard key gives every anagram in a
group one identical representation, so one hash lookup finds its bucket.

\subsection{Detailed Grouping Walkthrough}
\begin{dryrunbox}{Worked example}
\begin{center}
\begin{tabular}{@{}c c c@{}}
\toprule
Word & Sorted key & Bucket after insertion \\
\midrule
\texttt{eat} & \texttt{aet} & \texttt{[eat]} \\
\texttt{tea} & \texttt{aet} & \texttt{[eat,tea]} \\
\texttt{tan} & \texttt{ant} & \texttt{[tan]} \\
\texttt{ate} & \texttt{aet} & \texttt{[eat,tea,ate]} \\
\texttt{nat} & \texttt{ant} & \texttt{[tan,nat]} \\
\texttt{bat} & \texttt{abt} & \texttt{[bat]} \\
\bottomrule
\end{tabular}
\end{center}
\end{dryrunbox}

\subsection{Code Walkthrough}
For each original word, the code makes one copy named \texttt{key}, sorts that
copy, and appends the untouched word to \texttt{groups[key]}. After all words
are classified, each map value is moved into the returned vector. Moving
avoids copying every string bucket again.

\subsection{Complexity Details, Edge Cases, and Alternative}
If string lengths differ, a more precise sorting bound is
\(O(\sum_i k_i\log k_i)\). The map stores standard keys in addition to the
returned strings. Empty strings all use the empty key, and repeated identical
words remain repeated in the same group. Output order is unspecified because
\texttt{unordered\_map} iteration order is unspecified.

For lowercase English letters, a 26-element frequency signature avoids
sorting. It runs in \(O(\sum_i k_i)\) time but needs a collision-free way to
serialize counts, such as separators between decimal values.

\begin{visualbox}{Count the letters instead of sorting}
\centering
\begin{tikzpicture}[font=\small,>={Stealth[length=2.2mm,width=1.55mm]}]
  \node[dsnode] (word) {\texttt{tea}};
  \node[dsnode,right=17mm of word,fill=softorange] (count)
    {\(a:1,\ e:1,\ t:1\)};
  \node[dsnode,right=17mm of count,fill=softteal] (bucket)
    {same group as\\\texttt{eat}, \texttt{ate}};
  \draw[flowarrow,orange] (word)--node[edgelabel,above,font=\scriptsize]{count letters} (count);
  \draw[flowarrow,teal] (count)--node[edgelabel,above,font=\scriptsize]{use these counts} (bucket);
\end{tikzpicture}
\end{visualbox}
\begin{revisionbox}{How to explain it simply}
Anagrams become equal after sorting. I use that sorted form as a hash key and
append the original word to its bucket. This converts repeated pairwise
comparisons into one key construction and average constant-time lookup per
word.
\end{revisionbox}
\chaptercheckpoint
  {Items should be grouped by an equivalence relation.}
  {A standard string or frequency signature mapped to a vector of originals.}
  {Appending the transformed key instead of the original word.}
  {$O(nk\log k)$ with sorting, or $O(nk)$ with fixed alphabet counts.}
\judgeclassnote


\patternintoc{House Robber DP}
\problemstart{House Robber}{198}{Medium}{Rolling dynamic programming}{Two running values}

\subsection{The Problem}
Houses lie on a line. Return the maximum amount that can be robbed without
choosing two adjacent houses.

\subsection{From DP Array to Two Variables}
At each house, choose between skipping it and taking it with the best result
from two houses earlier: \(current=\max(previous, twoBack+money).\)
Only the previous two states are needed, which is the space-optimized approach
used in the displayed implementation.

\begin{invariantbox}{the two variables are the last two best prefixes}
Before processing a house, \texttt{previous} is the optimum for all houses
already seen and \texttt{twoBack} is the optimum before the immediately
previous house. Computing the new value before shifting keeps both choices
required by the update rule.
\end{invariantbox}

\begin{visualbox}{Compare skipping this house with taking it}
\centering
\begin{tikzpicture}[font=\small,>={Stealth[length=2.2mm,width=1.55mm]}]
  \node[dsnode,fill=softblue] (skip) {skip it\\keep 11};
  \node[dsnode,fill=softorange,below=8mm of skip] (take) {take it\\\(7+3=10\)};
  \node[dsnode,active,right=28mm of skip,yshift=-6mm] (cur) {best now\\11};
  \draw[flowarrow,navy] (skip.east)--(cur.west);
  \draw[flowarrow,orange] (take.east)--(cur.west);
  \node[left=12mm of skip,align=right,text=navy]
    {best before\\this house};
  \node[left=12mm of take,align=right,text=orange]
    {best two houses back\\plus this house};
\end{tikzpicture}

\smallskip
After the decision, \texttt{previous} becomes the new answer and the old
\texttt{previous} shifts into \texttt{twoBack} for the next house.
\end{visualbox}

\begin{dryrunbox}{Rolling states for \texttt{[2,7,9,3,1]}}
\begin{center}
\begin{tabular}{c|ccccc}
House value & 2 & 7 & 9 & 3 & 1 \\ \hline
Best so far & 2 & 7 & 11 & 11 & 12
\end{tabular}
\end{center}
\end{dryrunbox}

\subsection{Pseudocode}
\begin{lstlisting}[style=pseudocode]
twoBack = 0
previous = 0
for money in houses:
    current = max(previous, twoBack + money)
    twoBack = previous
    previous = current
return previous
\end{lstlisting}

\subsection{C++ Implementation}
\begin{lstlisting}[style=compactcode,caption={Complete C++ solution: House Robber}]
#include <algorithm>
#include <vector>

using namespace std;

class Solution {
public:
    int rob(vector<int>& nums) {
        int twoBack = 0;
        int previous = 0;
        for (const int money : nums) {
            const int current = max(previous, twoBack + money);
            twoBack = previous;
            previous = current;
        }
        return previous;
    }
};
\end{lstlisting}

\subsection{Why the Solution Works, Complexity, and Revision}
The update rule considers both possibilities for every best prefix: the last
house is skipped or included. Updating \texttt{twoBack} only after computing
\texttt{current} keeps the old states. Time is \(O(n)\), extra space is
\(O(1)\), and empty input naturally returns zero.

\begin{revisionbox}{House Robber pattern}
\textbf{State:} best answer for a prefix.\quad
\textbf{Choice:} skip current or take current plus the answer two positions
back.\quad \textbf{Signal:} adjacent choices cannot both be selected.
\end{revisionbox}

\subsection{From Subsets to Dynamic Programming}
Brute force chooses or skips every house, producing up to \(2^n\) subsets and
then rejecting those with adjacent selections. The useful fact is
that the best decision for the first \(i\) houses depends only on the best
answers for the first \(i-1\) and \(i-2\) houses.

For the current house, every best solution is in exactly one group:
\begin{itemize}
  \item skip it and keep the best previous result;
  \item take it, which forces the previous house to be skipped, and add its
  money to the best result two positions back.
\end{itemize}

\subsection{Detailed Rolling-State Trace}
\begin{dryrunbox}{Houses \texttt{[2,7,9,3,1]}}
\begin{center}
\begin{tabular}{@{}c c c c c@{}}
\toprule
Money & Two back & Previous & Take current & New best \\
\midrule
2 & 0 & 0 & 2 & 2 \\
7 & 0 & 2 & 7 & 7 \\
9 & 2 & 7 & 11 & 11 \\
3 & 7 & 11 & 10 & 11 \\
1 & 11 & 11 & 12 & 12 \\
\bottomrule
\end{tabular}
\end{center}
\end{dryrunbox}

\subsection{Code Walkthrough}
Before each iteration, \texttt{previous} is the best answer for the houses
already processed and \texttt{twoBack} is the answer one prefix earlier. The
code calculates \texttt{current} before overwriting either old value. It then
shifts the states forward. This keeps the update rule without allocating a
full DP array.

\subsection{Edge Cases and Interview Variations}
Empty input returns zero, one house returns its value, and two houses return
the larger value. The standard problem uses non-negative money. If negative
values were allowed and robbing no house remained valid, the zero
initialization would correctly ignore them. Variations include circular
houses, reconstructing which houses were chosen, or a minimum gap larger than
one between selected positions.

\begin{revisionbox}{How to explain it simply}
At each house I either skip it and keep the previous optimum, or take it
and add its money to the optimum from two houses back. Only those two earlier
states are needed, so the solution is \(O(n)\) time and \(O(1)\) extra space.
\end{revisionbox}
\chaptercheckpoint
  {Adjacent choices conflict and every prefix has a take-or-skip decision.}
  {The best values for the previous prefix and the prefix two steps back.}
  {Overwriting an old state before calculating the current decision.}
  {$O(n)$ time and $O(1)$ extra space.}
\judgeclassnote


\problemstart{House Robber II}{213}{Medium}{Two linear DP scans}{Dynamic programming}

\subsection{The Problem}
Houses form a circle, so the first and last houses are adjacent. Maximize the
money robbed without choosing adjacent houses.

\begin{conceptbox}{From a line to a circle}
The circular problem uses the same take-or-skip update rule as a linear row of
houses. Its only additional restriction is that the first and last houses
cannot both be selected.
\end{conceptbox}

\subsection{Break the Circle}
No valid solution can contain both the first and last houses. So, solve
two ordinary linear cases and take the larger result:
\begin{enumerate}
  \item houses \(0\) through \(n-2\), excluding the last;
  \item houses \(1\) through \(n-1\), excluding the first.
\end{enumerate}

\begin{visualbox}{Break the circle into two straight rows}
\centering
\begin{tikzpicture}[>={Stealth[length=2.2mm,width=1.55mm]},font=\small]
  \foreach \angle/\name/\v in {90/a/2,18/b/3,-54/c/2,-126/d/3,162/e/5} {
    \node[trienode] (\name) at (\angle:1.7cm) {\v};
  }
  \draw[-,navy,thick] (a)--(b)--(c)--(d)--(e)--(a);
  \node[right=18mm of c,align=left] {
    row A: skip the first house\\
    row B: skip the last house};
\end{tikzpicture}
\end{visualbox}

\subsection{Linear DP Update Rule}
For a range beginning at \(s\), let \(dp[i]\) be the best value using its first
\(i\) houses. Then \(dp[i]=\max(dp[i-1],dp[i-2]+\text{current house}).\)

\begin{visualbox}{Solve each straight row separately}
\centering
\begin{tikzpicture}[font=\small,>={Stealth[length=2.2mm,width=1.55mm]}]
  \node[anchor=east,text=navy,font=\bfseries] at (-0.5,1.35) {exclude last};
  \foreach \x/\v in {0/1,1/2,2/3} {
    \node[arraycell,minimum width=11mm] (ha\x) at (1.2*\x,1.35) {\v};
  }
  \node[arraycell,minimum width=11mm,fill=softteal] (da0) at (4.7,1.35) {1};
  \node[arraycell,minimum width=11mm,fill=softteal] (da1) at (5.9,1.35) {2};
  \node[arraycell,minimum width=11mm,fill=softteal] (da2) at (7.1,1.35) {4};
  \node[text=teal,font=\bfseries] at (3.7,1.35) {\(\rightarrow\) best totals};

  \node[anchor=east,text=navy,font=\bfseries] at (-0.5,-0.15) {exclude first};
  \foreach \x/\v in {0/2,1/3,2/1} {
    \node[arraycell,minimum width=11mm] (hb\x) at (1.2*\x,-0.15) {\v};
  }
  \node[arraycell,minimum width=11mm,fill=softorange] (db0) at (4.7,-0.15) {2};
  \node[arraycell,minimum width=11mm,fill=softorange] (db1) at (5.9,-0.15) {3};
  \node[arraycell,minimum width=11mm,fill=softorange] (db2) at (7.1,-0.15) {3};
  \node[text=orange,font=\bfseries] at (3.7,-0.15) {\(\rightarrow\) best totals};
  \node[below=8mm of db1,draw=teal,fill=softteal,rounded corners]
    {take \(\max(4,3)=4\)};
\end{tikzpicture}
\end{visualbox}

\subsection{Pseudocode}
\begin{lstlisting}[style=pseudocode]
robRange(start, end):
    dp[0] = 0
    dp[1] = nums[start]
    for each later house in the range:
        dp[i] = max(dp[i - 1], dp[i - 2] + houseValue)
    return final dp value

if there is only one house: return its money
return max(robRange(0, n - 2), robRange(1, n - 1))
\end{lstlisting}

\subsection{C++ Implementation}
\begin{lstlisting}[style=compactcode,caption={Complete C++ solution: House Robber II}]
#include <algorithm>
#include <vector>

using namespace std;

class Solution {
    int robRange(const vector<int>& nums, int start, int end) {
        const int length = end - start + 1;
        vector<int> dp(length + 1, 0);
        dp[1] = nums[start];
        for (int i = 2; i <= length; ++i) {
            const int value = nums[start + i - 1];
            dp[i] = max(dp[i - 1], dp[i - 2] + value);
        }
        return dp[length];
    }

public:
    int rob(vector<int>& nums) {
        const int n = static_cast<int>(nums.size());
        if (n == 0) return 0;
        if (n == 1) return nums[0];
        return max(robRange(nums, 0, n - 2),
                        robRange(nums, 1, n - 1));
    }
};
\end{lstlisting}

\subsection{Complexity and Space-Optimized Follow-up}
Both scans are linear, so time is \(O(n)\). The DP arrays use
\(O(n)\) space. Each scan can instead keep only the previous two states,
reducing extra space to \(O(1)\).

\begin{revisionbox}{How to explain it in an interview}
Because the first and last houses conflict, I solve two linear robber problems:
exclude the first or exclude the last. The linear update rule chooses between
skipping the current house and taking it with the best value two positions
back.
\end{revisionbox}

\subsection{Why Two Linear Cases Are Complete}
No legal robbery can contain both endpoint houses because they are adjacent
in the circle. Every legal solution belongs to at least one of two
sets: solutions that exclude the first house or solutions that exclude the
last. Taking the better optimum from those two sets covers every possibility.

\begin{invariantbox}{each helper call solves one complete legal linear case}
Within a chosen range, the DP value after each house is the best legal robbery
of that processed prefix. The two ranges together cover every circular
solution because no legal solution can contain both endpoint houses.
\end{invariantbox}

\subsection{Detailed Visual Dry Run}
\begin{dryrunbox}{Houses \texttt{[1,2,3,1]}}
\begin{center}
\begin{tabularx}{\linewidth}{@{}c c c X@{}}
\toprule
Linear case & Range & DP values & Best choice \\
\midrule
Exclude last & \texttt{[1,2,3]} & \(1,2,4\) & rob houses with 1 and 3 \\
Exclude first & \texttt{[2,3,1]} & \(2,3,3\) & rob the middle value 3 \\
Final & -- & \(\max(4,3)\) & return 4 \\
\bottomrule
\end{tabularx}
\end{center}
For \texttt{[2,3,2]}, both endpoint values 2 cannot be taken together, so the
middle house gives the answer 3.
\end{dryrunbox}

\subsection{Code Walkthrough}
\texttt{robRange} creates a DP table for one inclusive interval. The first
entry is the money in the first allowed house. Every later entry chooses
between the previous optimum and the current house plus the optimum two
positions earlier. The public method handles zero and one house before calling
the helper on the two non-empty ranges.

\subsection{Edge Cases, Space Optimization, and Follow-ups}
The special case for one house prevents both ranges from becoming empty. Two
houses return the larger value. Zero-valued houses work naturally. Each range
can use the two-variable rolling state from House Robber, reducing extra
space from \(O(n)\) to \(O(1)\) while preserving \(O(n)\) time.

\begin{revisionbox}{Practice checkpoint}
Write down the two ranges for five houses. Prove that a legal solution that
chooses neither endpoint is still considered by both ranges and so is
not lost.
\end{revisionbox}
\chaptercheckpoint
  {A linear adjacency rule becomes circular because the endpoints conflict.}
  {Two linear take-or-skip DP scans over complementary ranges.}
  {Running one DP over all houses and accidentally allowing both endpoints.}
  {$O(n)$ time and $O(n)$ DP-array space, reducible to $O(1)$.}
\judgeclassnote